\documentclass[12pt]{article}
\usepackage[utf8]{inputenc}
\usepackage[T1]{fontenc}
\usepackage[margin=1in]{geometry}
\usepackage{setspace}
\usepackage{xcolor}
\usepackage[colorlinks=true, urlcolor=blue!60!black, linkcolor=blue!60!black]{hyperref}
\setlength{\parindent}{2em}
\setlength{\parskip}{1.25em}
\renewcommand{\baselinestretch}{1.0}
\begin{document}
\begin{center}
{\LARGE \textit{Katzenbach v. Morgan}, 384 U.S. 641 (1966)\par}
\vspace{0.8em}
{\large Case Brief\par}
\end{center}
\vspace{1em}

\section{Memory Jogger}

In \href{https://www.courtlistener.com/opinion/107256/katzenbach-v-morgan/}{\textit{Katzenbach v. Morgan}}, the Supreme Court held that Section 5 of the Fourteenth Amendment is a positive grant of legislative power governed by \textit{McCulloch v. Maryland}, allowing Congress to override New York's English literacy test for Puerto Rican voters without a prior judicial finding that the test itself violated the Equal Protection Clause.

\section{Detailed Case Facts}

Article II, \S 1 of the New York Constitution (adopted in 1921) and \S\S 150 and 168 of the New York Election Law required new voters to demonstrate an ability to read and write English, either by passing a Regents literacy test or by presenting proof of completing sixth grade in an English-language school (\href{https://supreme.justia.com/cases/federal/us/384/641/}{384 U.S. at 644 \& n. 2}). Because Puerto Rican public and accredited private schools conducted instruction predominantly in Spanish, New York's English literacy requirement disenfranchised large numbers of the several hundred thousand United States citizens who had migrated from the Commonwealth of Puerto Rico to New York City (\href{https://supreme.justia.com/cases/federal/us/384/641/}{384 U.S. at 644--645 \& n. 3}).

During Senate floor debate on the Voting Rights Act of 1965, New York Senators Robert F. Kennedy and Jacob Javits sponsored an amendment---enacted without committee hearings as \S 4(e) of the Act, 79 Stat. 439, 42 U.S.C. \S 1973b(e) (\href{https://supreme.justia.com/cases/federal/us/384/641/}{384 U.S. at 643 n. 1, 645 n. 3, 669 n. 2/9}). Section 4(e) expressly declared that, ``to secure the rights under the fourteenth amendment of persons educated in American-flag schools in which the predominant classroom language was other than English,'' no person who successfully completed the sixth primary grade in a public school or accredited private school in Puerto Rico (or another United States jurisdiction) could be denied the right to vote in any federal, state, or local election because of an inability to read or write English (\href{https://supreme.justia.com/cases/federal/us/384/641/}{384 U.S. at 643 \& n. 1}).

Previously, in \textit{Lassiter v. Northampton County Board of Elections}, 360 U.S. 45 (1959), the Supreme Court had unanimously held that a facially neutral English literacy requirement did not, in all circumstances, violate the Fourteenth or Fifteenth Amendments (\href{https://supreme.justia.com/cases/federal/us/384/641/}{384 U.S. at 649}). Registered voters in New York City (John P. Morgan et al.) challenged \S 4(e), arguing that because New York's literacy test had not been adjudged unconstitutional under \textit{Lassiter}, Congress exceeded its \S 5 enforcement power and usurped powers reserved to the states under the Tenth Amendment (\href{https://supreme.justia.com/cases/federal/us/384/641/}{384 U.S. at 643--646}).

\section{Procedural History}

Pursuant to \S 14(b) of the Voting Rights Act of 1965, which vested exclusive jurisdiction to enjoin enforcement of the Act in the United States District Court for the District of Columbia, the New York City voter-appellees sued Attorney General Nicholas Katzenbach and later joined the New York City Board of Elections (which announced it would comply with \S 4(e) and defended the statute's validity), while the United States intervened as a defendant and the Attorney General of New York participated as \textit{amicus curiae} attacking \S 4(e) (\href{https://supreme.justia.com/cases/federal/us/384/641/}{384 U.S. at 645--646 \& n. 4}).

A three-judge District Court was convened under 28 U.S.C. \S\S 2282 and 2284. On cross-motions for summary judgment, the District Court held by a 2--1 vote (Circuit Judge McGowan dissenting) that Congress exceeded its constitutional powers and usurped state authority under the Tenth Amendment, and granted declaratory and injunctive relief against enforcing \S 4(e) (\textit{Morgan v. Katzenbach}, 247 F. Supp. 196 (D.D.C. 1965); \href{https://supreme.justia.com/cases/federal/us/384/641/}{384 U.S. at 646 \& n. 5}).

Both the Attorney General (No. 847) and the New York City Board of Elections (No. 877) took direct appeals to the Supreme Court under 28 U.S.C. \S\S 1252 and 1253; the Court noted probable jurisdiction at 382 U.S. 1007 (\href{https://supreme.justia.com/cases/federal/us/384/641/}{384 U.S. at 646, 659}). Argued April 18, 1966, alongside \textit{Cardona v. Power}, 384 U.S. 672 (1966), the case was decided June 13, 1966, reversing the District Court (\href{https://supreme.justia.com/cases/federal/us/384/641/}{384 U.S. at 641, 658}).

\section{Judicial Votes}

The Supreme Court reversed the District Court by a vote of \textbf{7--2} (\href{https://supreme.justia.com/cases/federal/us/384/641/}{384 U.S. at 643, 658--659}):

\begin{itemize}
\item \textbf{Majority Opinion:} Justice William J. Brennan, Jr. authored the opinion of the Court, joined in full by Chief Justice Earl Warren and Justices Hugo L. Black, Tom C. Clark, Byron R. White, and Abe Fortas, and joined in part by Justice William O. Douglas (\href{https://supreme.justia.com/cases/federal/us/384/641/}{384 U.S. at 643, 658--659}).
\item \textbf{Separate Statement / Partial Concurrence:} Justice William O. Douglas joined the Court's opinion except for the discussion at 384 U.S. at 656--658 rejecting the Fifth Amendment underinclusiveness challenge to the ``American-flag schools'' limitation, reserving judgment until a member of the excluded class presented that claim (\href{https://supreme.justia.com/cases/federal/us/384/641/}{384 U.S. at 658--659}).
\item \textbf{Dissenting Opinion:} Justice John Marshall Harlan II authored a dissenting opinion, joined by Justice Potter Stewart, applying to both \textit{Morgan} and \textit{Cardona v. Power} (\href{https://supreme.justia.com/cases/federal/us/384/641/}{384 U.S. at 659--671}).
\end{itemize}

\section{Holding}

\textbf{Judgment:} Reversed. Section 4(e) of the Voting Rights Act of 1965 is a valid exercise of Congress's enforcement power under \S 5 of the Fourteenth Amendment, and under the Supremacy Clause (U.S. Const. art. VI), New York's English literacy requirement cannot be enforced to the extent it conflicts with \S 4(e) (\href{https://supreme.justia.com/cases/federal/us/384/641/}{384 U.S. at 646--647, 658}).

\textbf{Rule Announced:} Congress's authority under \S 5 of the Fourteenth Amendment is not limited to abrogating state laws that the judiciary has already adjudged unconstitutional under \S 1 (\href{https://supreme.justia.com/cases/federal/us/384/641/}{384 U.S. at 648--650}). Rather, \S 5 is a positive grant of legislative power measured by the Necessary and Proper Clause standard of \textit{McCulloch v. Maryland}, 4 Wheat. 316, 421 (1819): legislation is ``appropriate'' if the Court can perceive a rational basis upon which Congress could conclude that the enactment enforces the Equal Protection Clause, is plainly adapted to that end, and is consistent with the letter and spirit of the Constitution (\href{https://supreme.justia.com/cases/federal/us/384/641/}{384 U.S. at 650--656}).

\textbf{Holding vs. Dicta:} The Court's holding is that \S 4(e) validly preempted New York's English literacy test as applied to sixth-grade graduates of Puerto Rican schools, and that \S 4(e)'s limitation to American-flag schools did not violate the Fifth Amendment because legislative reform ``may take one step at a time'' (\href{https://supreme.justia.com/cases/federal/us/384/641/}{384 U.S. at 656--658}). Crucially, Justice Brennan offered \textit{two} alternative rationales: (1) a \textit{remedial/prophylactic} rationale (enfranchising Puerto Ricans to help them secure nondiscriminatory public services such as schools, housing, and policing) and (2) a \textit{substantive/interpretive} rationale (allowing Congress itself to weigh competing considerations and determine that New York's literacy test worked an invidious discrimination in voting) (\href{https://supreme.justia.com/cases/federal/us/384/641/}{384 U.S. at 652--656}). In Footnote 10, Justice Brennan added what became known as the ``one-way ratchet'' theory---asserting that \S 5 grants Congress power only to expand, never to ``restrict, abrogate, or dilute,'' judicially recognized Fourteenth Amendment guarantees (\href{https://supreme.justia.com/cases/federal/us/384/641/}{384 U.S. at 651 n. 10}). That substantive-ratchet reading was later disavowed in \href{https://supreme.justia.com/cases/federal/us/521/507/}{\textit{City of Boerne v. Flores}}, 521 U.S. 507 (1997), which reframed \textit{Morgan} strictly around its remedial rationale and replaced \textit{McCulloch} deference with the ``congruence and proportionality'' test (\href{https://law.yale.edu/sites/default/files/documents/pdf/Faculty/Siegel_ProtectingTheConstitutionFromThePeople.pdf}{Post \& Siegel, 78 Ind. L.J. at 1--38}).

\section{Analysis of Opinions}

\subsection{Majority Opinion (Justice Brennan)}

Justice Brennan began by acknowledging that states generally set voter qualifications, but stressed that state voting rules remain fully subject to the Fourteenth Amendment (\href{https://supreme.justia.com/cases/federal/us/384/641/}{384 U.S. at 647}). He flatly rejected the New York Attorney General's contention that \S 5 permits Congress to invalidate a state law only after the judiciary declares that law unconstitutional under \S 1; such a view would reduce Congress to ``the insignificant role of abrogating only those state laws that the judicial branch was prepared to adjudge unconstitutional'' (\href{https://supreme.justia.com/cases/federal/us/384/641/}{384 U.S. at 648--649}). Accordingly, \textit{Lassiter} was inapposite because it addressed judicial invalidation under \S 1, not congressional legislation under \S 5 (\href{https://supreme.justia.com/cases/federal/us/384/641/}{384 U.S. at 649}).

Drawing on \textit{Ex parte Virginia}, 100 U.S. 339, 345--346 (1880), and \href{https://supreme.justia.com/cases/federal/us/383/301/}{\textit{South Carolina v. Katzenbach}}, 383 U.S. 301, 326 (1966), Justice Brennan held that \S 5 confers the same broad discretion as the Necessary and Proper Clause under \textit{McCulloch v. Maryland} (\href{https://supreme.justia.com/cases/federal/us/384/641/}{384 U.S. at 650--651}). Applying \textit{McCulloch}, he upheld \S 4(e) on two independent grounds. First, under the \textit{remedial} rationale, Congress could reasonably conclude that granting Puerto Rican residents the ballot---the right ``preservative of all rights''---would enable them to secure nondiscriminatory treatment in public schools, public housing, and law enforcement (\href{https://supreme.justia.com/cases/federal/us/384/641/}{384 U.S. at 652--653}). Second, under the \textit{substantive} rationale, Congress could independently weigh New York's asserted interests against evidence of nativist prejudice in the 1915--1921 adoption of the literacy test, the availability of Spanish-language newspapers and radio in New York City, and alternative ballot-assistance methods, and conclude that denying the vote to Spanish-literate Puerto Ricans was invidious discrimination (\href{https://supreme.justia.com/cases/federal/us/384/641/}{384 U.S. at 653--656 \& nn. 14--16}). In both inquiries, ``[i]t is enough that we be able to perceive a basis'' for Congress's judgment (\href{https://supreme.justia.com/cases/federal/us/384/641/}{384 U.S. at 653, 656}).

To answer Justice Harlan's warning that this logic would let Congress curtail constitutional rights, Justice Brennan added Footnote 10, insisting that \S 5 is strictly one-directional: Congress has ``no power to restrict, abrogate, or dilute these guarantees'' (\href{https://supreme.justia.com/cases/federal/us/384/641/}{384 U.S. at 651 n. 10}). Finally, he rejected the Fifth Amendment claim that \S 4(e) invidiously discriminated against citizens educated in foreign (non-American-flag) schools, citing \textit{Williamson v. Lee Optical Co.}, 348 U.S. 483, 489 (1955), for the rule that a reform measure expanding rights need not ``strike at all evils at the same time'' (\href{https://supreme.justia.com/cases/federal/us/384/641/}{384 U.S. at 656--658}).

\subsection{Partial Concurrence (Justice Douglas)}

Justice Douglas joined the Court's opinion on the core \S 5 holding, but declined to join the final section (384 U.S. at 656--658) upholding the American-flag school limitation against Fifth Amendment attack, reserving judgment on that issue until raised by a litigant actually excluded by that classification (\href{https://supreme.justia.com/cases/federal/us/384/641/}{384 U.S. at 658--659}).

\subsection{Dissenting Opinion (Justice Harlan, joined by Justice Stewart)}

Justice Harlan warned that upholding \S 4(e) sacrificed two constitutional fundamentals: ``the separation between the legislative and judicial function and the boundaries between federal and state political authority'' (\href{https://supreme.justia.com/cases/federal/us/384/641/}{384 U.S. at 659}). Addressing \textit{Cardona v. Power} first, he argued under \textit{Lassiter} that New York's fairly administered English literacy test satisfied rational-basis equal protection because government business, ballots, and statewide campaigns were conducted in English (\href{https://supreme.justia.com/cases/federal/us/384/641/}{384 U.S. at 660--664}).

Turning to \textit{Morgan}, Justice Harlan drew a sharp distinction between \textit{remedial} and \textit{substantive} power: while \S 5 grants Congress wide remedial latitude once a constitutional violation is established (as in \textit{Ex parte Virginia} and \href{https://supreme.justia.com/cases/federal/us/383/301/}{\textit{South Carolina v. Katzenbach}}), ``it is a judicial question whether the condition with which Congress has thus sought to deal is, in truth, an infringement of the Constitution'' (\href{https://supreme.justia.com/cases/federal/us/384/641/}{384 U.S. at 666--667}). Allowing Congress to redefine the substantive scope of \S 1 under a mere rationality test would permit Congress to override Supreme Court precedent---and, logically, to \textit{dilute} equal protection and due process decisions just as easily as to expand them (\href{https://supreme.justia.com/cases/federal/us/384/641/}{384 U.S. at 667--668}). Moreover, unlike \href{https://supreme.justia.com/cases/federal/us/379/241/}{\textit{Heart of Atlanta Motel}}, \href{https://supreme.justia.com/cases/federal/us/379/294/}{\textit{McClung}}, and \href{https://supreme.justia.com/cases/federal/us/383/301/}{\textit{South Carolina v. Katzenbach}}, Congress compiled \textit{no} legislative record or committee hearings documenting discrimination against Puerto Ricans in New York's public services (\href{https://supreme.justia.com/cases/federal/us/384/641/}{384 U.S. at 668--669 \& n. 2/9}). Overriding a presumptively valid state voting law by ``mere \textit{ipse dixit}'' threatened to let the Fourteenth Amendment swallow primary state legislative competence (\href{https://supreme.justia.com/cases/federal/us/384/641/}{384 U.S. at 670--671}).

\section{Examples: Future Applications}

\begin{enumerate}
\item \textbf{Same-Side (Valid under \textit{Morgan} and Post-\textit{Boerne} Doctrine):} Following extensive legislative hearings documenting state courtrooms' systematic exclusion of wheelchair users from judicial proceedings, Congress enacts a statute requiring states to provide physical accessibility to courthouses and abrogating state sovereign immunity for violations. \textit{Why it falls on the same side:} Just as \textit{Morgan}'s remedial theory upheld \S 4(e) to protect fundamental rights, the Court in \href{https://supreme.justia.com/cases/federal/us/541/509/}{\textit{Tennessee v. Lane}}, 541 U.S. 509 (2004), upheld Title II of the Americans with Disabilities Act as applied to courthouse access because it targeted a documented history of unconstitutional denials of a fundamental right (access to courts) with a congruent and proportional remedy.

\item \textbf{Same-Side (Valid under \textit{Morgan}'s Remedial Rationale):} Congress enacts a nationwide ban on state literacy tests for voting after compiling a substantial legislative record showing that literacy tests perpetuate the effects of past racially segregated public education across the country. \textit{Why it falls on the same side:} Even Justices skeptical of \textit{Morgan}'s substantive ratchet theory unanimously upheld the nationwide literacy-test ban in \href{https://supreme.justia.com/cases/federal/us/400/112/}{\textit{Oregon v. Mitchell}}, 400 U.S. 112 (1970), as an appropriate remedial measure enforcing the Fourteenth and Fifteenth Amendments against documented racial discrimination.

\item \textbf{Opposite-Side (Invalid under \textit{City of Boerne v. Flores}):} Disagreeing with a Supreme Court Free Exercise ruling that neutral, generally applicable state laws do not trigger strict scrutiny, Congress passes a nationwide statute under \S 5 imposing strict scrutiny on every state and local law that substantially burdens religious exercise, without a record of widespread state religious bigotry. \textit{Why it falls on the opposite side:} In \href{https://supreme.justia.com/cases/federal/us/521/507/}{\textit{City of Boerne v. Flores}}, 521 U.S. 507 (1997), the Court struck down the Religious Freedom Restoration Act as applied to the states, adopting Justice Harlan's \textit{Morgan} dissent and holding that Congress may not use \S 5 to enact substantive changes to constitutional rights that lack congruence and proportionality to documented \S 1 violations.

\item \textbf{Opposite-Side (Invalid under \textit{Morgan} Footnote 10 and \textit{Boerne}):} Congress invokes \S 5 of the Fourteenth Amendment to enact a statute declaring that state laws requiring racially segregated public schools do not violate the Equal Protection Clause so long as facilities receive equal funding. \textit{Why it falls on the opposite side:} Even under \textit{Morgan}'s broadest reading, Footnote 10 explicitly barred Congress from using \S 5 to ``restrict, abrogate, or dilute'' judicially established Fourteenth Amendment guarantees (\href{https://supreme.justia.com/cases/federal/us/384/641/}{384 U.S. at 651 n. 10}), and \href{https://supreme.justia.com/cases/federal/us/521/507/}{\textit{City of Boerne}} equally forbids congressional alteration of \S 1 rights.

\item \textbf{Fence-Sitter (Borderline Application):} Congress enacts a statute under \S 5 lowering the voting age to 18 in all state and local elections, reasoning that 18-to-20-year-olds are subject to the military draft and taxation and that excluding them from the franchise denies equal protection. \textit{Why it is a fence-sitter:} Under \textit{Morgan}'s substantive rational-basis theory, Congress could rationally weigh youth maturity against state interests and expand the franchise. Under Justice Harlan's dissent and later doctrine, however, age classifications receive rational-basis review and Article I, \S 2 reserves state voter qualifications to the states. Indeed, this exact issue split the Court 5--4 in \href{https://supreme.justia.com/cases/federal/us/400/112/}{\textit{Oregon v. Mitchell}}, 400 U.S. 112 (1970), which invalidated the 18-year-old vote for state elections (prompting the Twenty-Sixth Amendment) while upholding it for federal elections.
\end{enumerate}

\section{Critique}

\textbf{Originalist and Textualist Perspectives:} Originalist, textualist, and structural critics echo Justice Harlan's dissent in arguing that the text of \S 5---granting power to ``\textit{enforce}'' the provisions of the Fourteenth Amendment---presupposes a fixed substantive baseline in \S 1 rather than an open-ended license for Congress to redefine constitutional rights (\href{https://supreme.justia.com/cases/federal/us/521/507/}{\textit{City of Boerne v. Flores}}, 521 U.S. at 519--524; \href{https://scholarship.richmond.edu/cgi/viewcontent.cgi?article=2278&context=lawreview}{University of Richmond Law Review}). Stephen L. Carter demonstrated both that the framers of the Fourteenth Amendment rejected broader ``necessary and proper'' language in favor of ``enforce, by appropriate legislation,'' and that \textit{Morgan}'s ``one-way ratchet'' breaks down whenever expanding one group's asserted right (such as a proposed federal ``Human Life Bill'' defining fetal personhood under \S 5) collides with another person's judicially recognized right (such as reproductive liberty) (\href{https://chicagounbound.uchicago.edu/context/uclrev/article/4477/viewcontent/53UChiLRev819.pdf}{Carter, 53 U. Chi. L. Rev. at 819--828}; \href{https://research.ualr.edu/cgi/viewcontent.cgi?article=1624&context=bowen_lawreview}{UA Little Rock Law Review}). From this perspective, \textit{Morgan}'s substantive rationale bypasses the Article V amendment process and undermines \textit{Marbury v. Madison}'s allocation of constitutional interpretation to the judiciary (\href{https://scholarship.richmond.edu/cgi/viewcontent.cgi?article=2278&context=lawreview}{University of Richmond Law Review}; \href{https://constitutioncenter.org/the-constitution/amendments/amendment-xiv/clauses/703}{National Constitution Center}).

\textbf{Progressive and Doctrinal Perspectives:} Conversely, many progressive and constitutional-theory scholars defend \href{https://www.courtlistener.com/opinion/107256/katzenbach-v-morgan/}{\textit{Katzenbach v. Morgan}} as a vital model of \textit{popular constitutionalism} and shared interpretive responsibility. Robert C. Post and Reva B. Siegel argue that the Court's later retreat from \textit{Morgan} in \href{https://supreme.justia.com/cases/federal/us/521/507/}{\textit{City of Boerne}} erected ``juricentric'' restrictions that strip the democratic branches of their historic Reconstruction-era role in giving concrete meaning to the Fourteenth Amendment's broad guarantees (\href{https://law.yale.edu/sites/default/files/documents/pdf/Faculty/Siegel_ProtectingTheConstitutionFromThePeople.pdf}{Post \& Siegel, 78 Ind. L.J. at 1--38}; \href{https://law.yale.edu/sites/default/files/documents/pdf/Faculty/Siegel_PopularConstitutionalismDepartmentalismAndJudicialSupremacy.pdf}{Yale Law School}). Similarly, historical scholarship notes that the 39th Congress intended \S 5 primarily to empower \textit{Congress}, not exclusively the courts, to enforce Reconstruction ideals (\href{https://scholarlycommons.law.northwestern.edu/cgi/viewcontent.cgi?article=1349&context=nulr}{Northwestern University Law Review}). Yet even sympathetic doctrinalists---such as Archibald Cox---cautioned early on that applying \textit{McCulloch}'s deferential standard without a clear limiting principle risked leaving judicial checks on congressional \S 5 power virtually nonexistent (\href{https://research.ualr.edu/cgi/viewcontent.cgi?article=1624&context=bowen_lawreview}{UA Little Rock Law Review}).

\textbf{Underlying Competing Values:} \textit{Katzenbach v. Morgan} squarely pits \textit{judicial supremacy and federalism} (guarding the Court's role as final arbiter of constitutional meaning and protecting state authority over local governance and elections) against \textit{democratic constitutionalism and national civil rights enforcement} (empowering a nationally accountable Congress to use its superior factfinding and political responsiveness to realize equal citizenship).

\section{Key Quotations}

\begin{quote}
``Thus, the \textit{McCulloch v. Maryland} standard is the measure of what constitutes `appropriate legislation' under \S 5 of the Fourteenth Amendment. Correctly viewed, \S 5 is a positive grant of legislative power authorizing Congress to exercise its discretion in determining whether and what legislation is needed to secure the guarantees of the Fourteenth Amendment.''
\end{quote}
\noindent --- Justice William J. Brennan, Jr., Opinion of the Court, 384 U.S. at 651 (\href{https://supreme.justia.com/cases/federal/us/384/641/}{Justia} \textbar{} \href{https://www.courtlistener.com/opinion/107256/katzenbach-v-morgan/}{CourtListener}).

\begin{quote}
``It is not for us to review the congressional resolution of these factors. It is enough that we be able to perceive a basis upon which the Congress might resolve the conflict as it did.''
\end{quote}
\noindent --- Justice William J. Brennan, Jr., Opinion of the Court, 384 U.S. at 653 (\href{https://supreme.justia.com/cases/federal/us/384/641/}{Justia} \textbar{} \href{https://www.courtlistener.com/opinion/107256/katzenbach-v-morgan/}{CourtListener}).

\begin{quote}
``We emphasize that Congress' power under \S 5 is limited to adopting measures to enforce the guarantees of the Amendment; \S 5 grants Congress no power to restrict, abrogate, or dilute these guarantees.''
\end{quote}
\noindent --- Justice William J. Brennan, Jr., Opinion of the Court, 384 U.S. at 651 n. 10 (\href{https://supreme.justia.com/cases/federal/us/384/641/}{Justia} \textbar{} \href{https://www.courtlistener.com/opinion/107256/katzenbach-v-morgan/}{CourtListener}).

\begin{quote}
``Worthy as its purposes may be thought by many, I do not see how \S 4(e) of the Voting Rights Act of 1965, 79 Stat. 439, 42 U.S.C. \S 1973b(e) (1964 ed. Supp. I), can be sustained except at the sacrifice of fundamentals in the American constitutional system --- the separation between the legislative and judicial function and the boundaries between federal and state political authority.''
\end{quote}
\noindent --- Justice John Marshall Harlan II, Dissenting Opinion, 384 U.S. at 659 (\href{https://supreme.justia.com/cases/federal/us/384/641/}{Justia} \textbar{} \href{https://www.courtlistener.com/opinion/107256/katzenbach-v-morgan/}{CourtListener}).

\begin{quote}
``In effect, the Court reads \S 5 of the Fourteenth Amendment as giving Congress the power to define the substantive scope of the Amendment. If that indeed be the true reach of \S 5, then I do not see why Congress should not be able as well to exercise its \S 5 `discretion' by enacting statutes so as, in effect, to dilute equal protection and due process decisions of this Court.''
\end{quote}
\noindent --- Justice John Marshall Harlan II, Dissenting Opinion, 384 U.S. at 668 (\href{https://supreme.justia.com/cases/federal/us/384/641/}{Justia} \textbar{} \href{https://www.courtlistener.com/opinion/107256/katzenbach-v-morgan/}{CourtListener}).

\end{document}
